
\documentclass[12pt]{article}

\usepackage[utf8]{inputenc}
\usepackage[margin=1in]{geometry}
\usepackage{graphicx}
\usepackage{amsmath,amssymb}
\usepackage{tikz}
\usepackage[parfill]{parskip}
\usepackage{systeme}
\usepackage{longtable} 

% ---- Math shortcuts ----
\newcommand{\E}{\mathbb{E}}
\newcommand{\Prb}{\mathbb{P}}
\newcommand{\X}{\mathcal{X}}
\begin{document}
\title{Risk Estimation and Control in Dynamic Systems}
\author{Notes by Miguel Bueno}
\maketitle

\newpage

\section{Abstract}

We consider a population of entities $\mathcal{X}$ evolving over time $t$, where each entity $x_i$ has a latent true label $Y_i$ and engagement weight $E_{it}$. The target estimand is the engagement-weighted prevalence $\mu_E = \mathbb{E}_{x \sim p_t}[Y(x)] \mid  p_t(x) \propto E_{it}$, and the global objective is to minimize an upper confidence bound on a policy-dependent loss functional $\mathcal{L}(\tau,\mathcal{G})$ at $t+1$ by jointly optimizing the deployment policy $\tau$ and the labeling policy $\mathcal{G}$.

We have available a set of resources $\mathcal{R}$ categorized by type $\kappa = \rho(r)$.\footnote{E.g. deep neural-networks , large language models, novice and expert human annotators, etc.}. Each resource is characterized by a Resource Profile $\Phi_\kappa = \langle c_\kappa, \Omega_\kappa, \mathbf{N}_\kappa \rangle$, representing its marginal cost, processing capacity, and noise. \footnote{Such characteristics can evolve over time. However, we assume they are relatively stable in the short-run.} \footnote{In reality noise is often heteroscedastic, arising from resource capabilities and latent task complexity.} Under a fixed global budget $B_t$ and total system bandwidth $\Omega_t$, the policies $\tau$ and $\mathcal{G}$ are constrained such that aggregate cost and throughput satisfy budget and capacity constraints (e.g., $\sum_{\kappa} c_\kappa n_\kappa \leq B_t$ and $\sum_{\kappa} \Omega_\kappa \leq \Omega_t$). 

A resource type $\kappa$ is said to be high-fidelity if it exhibits low noise, $\mathbf{N}_{\kappa} \approx \mathbf{I}$, but is constrained in cost and/or capacity, such that $c_{\kappa}$ is large and $\Omega_\kappa \ll |\mathcal{X}|$. Conversely, a resource type $\kappa$ is low-fidelity if it is inexpensive and scalable, with $c_\kappa \ll c_{\mathrm{HF}}$ and $\Omega_\kappa \gtrsim |\mathcal{X}|$, but exhibits substantially higher noise \footnote{As reflected by $\operatorname{tr}(\mathbf{N}_\kappa) \gg \operatorname{tr}(\mathbf{N}_{\mathrm{HF}})$}. A resource type $\kappa$ is mid-fidelity if it lies between these extremes in noise, cost, and capacity, while still remaining meaningfully capacity-constrained relative to the low-fidelity resource.

High- and mid- fidelity resources are used to produce reference labels $Y_i^*$ for estimating $\mu_E$. We construct a sampled dataset $\mathcal{D}_{n,t} = \{X_i \mid R_{it}=1\} \subset \mathcal{X}$ with inclusion probabilities $\pi_{it}$. To efficiently estimate $\mu_E$, samples are drawn under an importance distribution $q(x)$ (potentially stratified over $\mathcal{H}$ and pooled over horizon $T$), with importance weights $w_i = p_{it} \cdot (q_{it})^{-1}$. Accounting for an availability probability $P(A_{i}=1|Y_{i}=1) \leq 1$ at the labeling stage, each available sampled entity $X_i$ is evaluated by a panel of mid-fidelity resources $\mathcal{P}_i \subset \mathcal{R}$ under policy $\mathcal{G}$. Each resource $r \in \mathcal{P}_i$ produces a valuation $v_{ir}$ \footnote{For human annotators, $v_{ir}$ is typically derived from structured ordinal responses $S_{ir\ell}$ over a template $\mathcal{Q}$.}, which is aggregated via $\mathcal{A}(\cdot)$ to produce an aggregated valuation $\hat{V}_{i} = \mathcal{A}(\{v_{ir}\}_{r \in \mathcal{P}_i})$ for all $x_{i} \in \mathcal{D}_{n,t}$. These valuations are mapped to discrete outcomes using a decision threshold $\eta$. The predicted label $\hat{Y}_{i} = \mathbb{I}(\hat{V}_{i}>\eta)$ is then used to estimate $\mu_{E}$.

To minimize redundant expenditure, we implement adaptive task assignment where panel size $|\mathcal{P}_i| = \phi(\alpha_i, \epsilon)$ is dynamically scaled.\footnote{Task difficulty is frequently proxied by agreement ($\alpha_i$) calculated over an initial rater subset.} When $\mathcal{P}_i$ used to construct $\hat{Y}_{i}$ consists of mid-fidelity resources, we occasionally leverage high-fidelity resources with difference estimators to reduce the variance of the estimator.\footnote{E.g. prediction-power-inference (PPI), double-sampling, etc.}

Where appropriate, low-fidelity resources are trained, fine-tuned, and tested using $Y_i^*$. They, in turn, generate valuations $v_{ir}$ for $\mathcal{X}$. These valuations drive $\tau$, inducing dynamics governed by $p_{t+1} = \mathcal{F}(p_t, \tau_t)$. After calibration via $g(\cdot)$, these valuations inform both $\pi_{it}$ for the production of $Y_i^*$ and the deployment policy $\tau$, which maps calibrated scores to exposure levels $a_\tau(x) \in [0,1]$ interpreted as a multiplicative ranking factor. \footnote{$a_\tau(x)=1$ corresponds to no demotion and $a_\tau(x)=0$ corresponds to complete filtering.}\footnote{Where explainability is required, $g(v_{ir})$ may be accompanied by a vector of interpretive features.}

Low-fidelity resource performance is assessed using a self-normalized, importance-weighted confusion matrix $\mathbf{M}(\eta)$, yielding weighted precision $p_{\omega m}(\mathbf{M})$ and recall $r_{\omega m}(\mathbf{M})$ and other metrics. Evaluation is conducted under a counterfactual no-policy baseline by weighting observations using predicted counterfactual engagement $\widehat{E}_{it}$ to avoid demotion bias. When $\mathcal{P}_i$ used to construct $Y_i^*$ consists of mid-fidelity resources, we refer to the reference label as noisy and denote it $\tilde{Y}_i$. To evaluate each panel, a subsample $\mathcal{D}_{n^{(2)}} \subset \mathcal{D}_{n,t}$ is adjudicated by high-fidelity resources. All resources are evaluated for fairness for protected partitions of entities, $\mathcal{M} \subset \mathcal{X}$ using the raw\footnote{Non-engagement weighted} false positive rate difference, $\Delta_{\mathcal{M},\mathcal{X}}\text{FPR}$, where the baseline represents the aggregate performance across the entire population $\mathcal{X}$. To optimize $B_{t}$, the system utilizes adaptive sampling schemes designed to maximize the informational utility of each acquired label.\footnote{E.g. multiple importance sampling (MIS), adaptive sampling, and sequential stopping rules.}

We can influence resource characteristics $\Phi_{\kappa}$ through treatments $W$ \footnote{E.g. Changes to a template $\mathcal{Q}$, refinements to the labeling policy $\mathcal{G}$, etc.} whose effects are often estimated via AB tests\footnote{i.e randomized control trials} that condition on additional features $Z_{i}$. Furthermore, we can identify resources $r$ with unusually unsatisfactory characteristics relative to those of the same type $\kappa$ and exclude them from $\mathcal{R}$.\footnote{When $r$ corresponds to a human annotator, such interventions must comply with applicable co-employment laws.}

Uncertainty in all estimators $\hat{\theta}$ is quantified via a resampling operator $\mathcal{B}$, which induces bootstrap replicates $\mathcal{D}^{(b)}_{n,t} \sim \mathcal{B}(\mathcal{D}_{n,t})$, with estimates $\hat{\theta}^{(b)}$ summarized through their empirical distribution.\footnote{$\mathcal{B}$ is typically leverages a paired pigeonhole bootstrapping but can leverage other bootstraps.} 

Since $|\mathcal{X}|$ and most $|\mathcal{D}_{n,t}|$ are large, engines are configured to use Poisson sampling variants to mimic multinomial draws and facilitate high-throughput.\footnote{These are co-deployed with Gumbel-Max-style weighted reservoir sampling procedures.}

Collectively, these components define a closed-loop system in which sampling, labeling, evaluation, and deployment interact through shared resource constraints and evolving population dynamics. The system objective is therefore to solve the constrained optimization problem $\min_{\tau, \mathcal{G}} \; \operatorname{UCB}\!\left(\mathcal{L}(\tau,\mathcal{G}) \right)$ subject to $\sum_{i \in \mathcal{D}_{n,t}} \sum_{r \in \mathcal{P}_i} c_{\rho(r)} \le B_t$ and $\sum_{i \in \mathcal{D}_{n,t}} \mathbf{1}\{\rho(r)=\kappa,\; r \in \mathcal{P}_i\} \le \Omega_\kappa,  \forall \kappa $.

\newpage
\section{Glossary}

\renewcommand{\arraystretch}{0.95}
\setlength{\tabcolsep}{4pt}

\begin{longtable}{p{0.32\linewidth} p{0.64\linewidth}}

\multicolumn{2}{l}{\textbf{Population and Target}} \\
$\mathcal{X}$ & Population of entities \\
$x, x_i$ & Entity in the population \\
$i$ & Entity index \\
$t$ & Time index \\
$Y_i$ & Latent true label (1 = positive class) \\
$E_{it}$ & Engagement weight of entity $x_i$ at time $t$ \\
$p_t(x)$ & Nominal (traffic-weighted) distribution, $p_t(x) \propto E_{it}$ \\
$\mu_E = \mathbb{E}_{x \sim p_t}[Y(x)]$ & Engagement-weighted prevalence \\
\\[3pt]

\multicolumn{2}{l}{\textbf{Sampling}} \\
$X_i$ & Sampled entity \\
$R_{it}$ & Sampling indicator \\
$\pi_{it} = \mathbb{P}(R_{it}=1)$ & Inclusion probability \\
$\mathcal{D}_{n,t}$ & Sampled dataset $\{X_i \mid R_{it}=1\}$ \\
$q(x)$ & Importance sampling distribution \\
$w_i = \frac{p_t(X_i)}{q(X_i)}$ & Importance weight \\
$\mathcal{H}$ & Set of strata \\
$h$ & Stratum index \\
$T$ & Time horizon for pooled sampling \\
\\[3pt]

\multicolumn{2}{l}{\textbf{Resources and Policies}} \\
$\mathcal{R}$ & Set of resources \\
$r$ & Resource index \\
$\mathcal{K}$ & Set of resource types \\
$\kappa$ & Resource type index \\
$\rho(r)$ & Mapping from resource to type \\
$\Phi_\kappa = \langle c_\kappa, \Omega_\kappa, \mathbf{N}_\kappa \rangle$ & Resource profile \\
$c_\kappa$ & Marginal cost \\
$\Omega_\kappa$ & Processing capacity \\
$\mathbf{N}_\kappa$ & Noise structure \\
$B_t$ & Total budget \\
$\mathcal{G}$ & Labeling policy \\
\\[3pt]

\multicolumn{2}{l}{\textbf{Annotation and Valuation}} \\
$\mathcal{P}_i \subset \mathcal{R}$ & Panel of resources \\
$v_{ir}$ & Valuation \\
$\mathcal{A}(\cdot)$ & Aggregation function \\
$\hat{V}_i = \mathcal{A}(\{v_{ir}\})$ & Aggregated valuation \\
$\eta$ & Decision threshold \\
$\hat{Y}_i = \mathbb{I}(\hat{V}_i > \eta)$ & Predicted label \\
$\tilde{Y}_i$ & Noisy reference label \\
$\mathcal{D}_{n^{(2)}}$ & High-fidelity subsample \\
\\[3pt]

\multicolumn{2}{l}{\textbf{Structured Human Annotation}} \\
$\mathcal{R}^{(H)}$ & Human annotators \\
$\mathcal{Q}$ & Annotation questions \\
$\ell$ & Question index \\
$S_{ir\ell}$ & Ordinal response \\
$v_{ir}$ & Derived valuation \\
\\[3pt]

\multicolumn{2}{l}{\textbf{Prediction and Deployment}} \\
$v(x)$ & Scalable valuation \\
$g(\cdot)$ & Calibration \\
$a_\tau(x) \in [0,1]$ & Exposure multiplier \\
$\tau$ & Deployment policy \\
$p_{t+1} = \mathcal{F}(p_t, \tau_t)$ & Population dynamics \\
\\[3pt]

\multicolumn{2}{l}{\textbf{Evaluation}} \\
$\mathbf{M}(\eta)$ & Weighted confusion matrix \\
$p_{\omega}(\mathbf{M})$ & Weighted precision \\
$r_{\omega}(\mathbf{M})$ & Weighted recall \\
$\hat{E}_{it}$ & Counterfactual engagement \\
$\mathcal{M}$ & Protected subset \\
$\Delta_{\mathcal{M},\mathcal{X}}\mathrm{FPR}$ & FPR difference \\
\\[3pt]

\multicolumn{2}{l}{\textbf{Objective and Optimization}} \\
$\mathcal{L}(\tau,\mathcal{G})$ & Loss functional \\
$\operatorname{UCB}(\cdot)$ & Upper confidence bound \\
$\min_{\tau,\mathcal{G}} \operatorname{UCB}(\mathcal{L})$ & System objective \\
$n_\kappa$ & Assignments per type \\
\\[3pt]

\multicolumn{2}{l}{\textbf{Uncertainty}} \\
$\hat{\theta}$ & Estimator \\
$\mathcal{B}$ & Bootstrap operator \\
$b$ & Bootstrap index \\
$\mathcal{D}^{(b)}_{n,t}$ & Resampled dataset \\
$\hat{\theta}^{(b)}$ & Bootstrap estimate \\
\\[3pt]

\multicolumn{2}{l}{\textbf{Causal Variables}} \\
$W$ & Treatment \\
$Z_i$ & Covariates \\

\end{longtable}

\newpage

\end{document}

